\section{The finite all-input word}
\label{app:finite-unit}

\subsection{Root lines, projectors and binary incidence}
Take one representative of each of the $120$ lines
$2e_i\pm2e_j$ ($i<j$) and $(1,\pm1,\ldots,\pm1)$ with even
sign parity. Each vector $v_s$ has squared norm eight. Distinct
representatives have inner products $0,\pm4$. Put
$P_s=v_sv_s^T/8$ over $\mathbb Q$, or reduce the same expression over
$\F_3$. Thus $P_s$ is a rank-one idempotent; orthogonal root lines
give mutually annihilating idempotents in both interpretations.

Let $A_{st}=1$ for distinct orthogonal lines and zero otherwise.
Over $\F_2$ put $B=A+I$. At scale $\alpha_s=v_s/2$ choose
\[
\begin{split}
\beta_1&=(1,-1,-1,-1,-1,-1,-1,1)/2,\qquad\beta_2=e_1+e_2,\\
\beta_3&=-e_1+e_2,\quad\beta_4=-e_2+e_3,\quad
\beta_5=-e_3+e_4,\\
\beta_6&=-e_4+e_5,\quad\beta_7=-e_5+e_6,\quad
\beta_8=-e_6+e_7.
\end{split}
\]
Their Gram matrix $S$ is the E8 tree matrix with diagonal two and
three arms of lengths $1,2,4$. A length-$l$ arm has determinant
$l+1$ and inverse endpoint entry $l/(l+1)$. Taking its Schur
complement gives
\[
 \det S=30(2-1/2-2/3-4/5)=1.
\]
The vectors $d_s=(\alpha_s\cdot\beta_i)_i$ are integral: two half
roots have even Hamming distance, and integer/half-root products are
integral. Since the $\beta_i$ form a basis,
$\alpha_s\cdot\alpha_t=d_s^TS^{-1}d_t$. The integral matrix
$S^{-1}$ therefore gives the explicit binary factors
\[
 R_s=(1,d_s^T),\qquad G_t=(1,S^{-1}d_t)^T,\qquad RG=B.
\]
The equality follows entrywise: the root inner product is even on
the diagonal and on orthogonal pairs, and odd on every other pair.

\subsection{Fifteen bases and the rectangle list}
Index the eight coordinates by $x\in\F_2^3$. For each nonzero $a$,
the four pairs $\{x,x+a\}$ give the eight orthogonal roots
$2e_x\pm2e_{x+a}$. These seven bases partition the $56$ sparse roots.
The binary simplex code $(a\cdot x)_{x\ne0}$ has eight words; its
nonzero words have weight four. Its eight cosets in the even-weight
length-seven code therefore give eight orthogonal bases of sign
vectors with coordinate zero fixed positive. Take the least unused
integer word as each coset representative. This partitions the
remaining $64$ root lines.

Give each basis a balanced tree on eight positions. Enumerate the
$21$ flags by choosing $a,b\in\{1,\ldots,7\}$ with $b\ne a$ and
$b<(b\mathbin{\mathrm{xor}}a)$, and then the least $c$ outside their
span. Relabel position $x$ by
$4(a\cdot x)+2(b\cdot x)+(c\cdot x)$. In the base order above use
the following zero-based source and target flag indices:
\[
\begin{split}
 f_X&=(9,0,0,4,0,16,1,9,15,9,9,1,13,1,15),\\
 f_Y&=(10,2,0,10,4,15,0,10,16,11,14,0,14,10,15).
\end{split}
\]
The accompanying file \path{checks/expected/frozen_rectangles.json}
lists the $3782$ rectangles. Each source and target set is a node
of the corresponding tree. The verification algorithm checks that
each rectangle consists entirely of orthogonal pairs and that its
directed pairs are disjoint from all earlier rectangles; finally it
checks that every directed orthogonal pair has been covered. All
sets have at most eight elements, so this is a finite exact
certificate with no asymptotic premise. Its histogram is the one in
Lemma~\ref{lem:unit}; $\sum abS_{a,b}=7560$.

Let $V_{e,t}=1$ when $t$ is in rectangle $e$'s source set and
$J_{s,e}=1$ when $s$ is in its target set. Then $JV=A$ and
\begin{equation}
 RG+JV=I\quad\text{over }\F_2.\label{eq:binary-incidence}
\end{equation}
This binary identity concerns payload additions. It is not reduced
from a rational identity.

\subsection{The common-frame program}
Write $\mathcal K_U$ for the paired phase associated with a projector
$P_U$. A field in frame $F$ holds $F$ times its logical value.
Initially assign data field $X_t$ frame $\mathcal K_{P_t}$ and
fields $Y_s,Z_e$ frame $I$. This assignment is notation: its logical
$X_t$ value is $\mathcal K_{P_t}^{-1}X_t$, and it requires no unpaid
initial query. The following chronological program uses independent
arbitrary $Y$ and side inputs $Z$.

\begin{enumerate}
\item Subtract $JZ$ from $Y$, all in frame $I$.
\item Visit source-tree nodes from leaves to root. Grow participating
$X_t$ fields to the node frame $\mathcal K_U$, advance every contacting
side to that same frame, and add its $V$-row sum of $X_t$. Each contact
therefore takes place in a common frame. At the root all $X$ fields
are in frame $\mathcal K_I$.
\item Form the nine $G$-row increments as owned fields. On each apply
$\mathcal K_I^{-1}$, implemented on disjoint rank-seven and rank-one
pieces, scatter it by $R$ to the still unadvanced $Y$, and erase it.
All inverse queries are paid.
\item Visit target-tree nodes from root to leaves. At a node of span
$W$, advance its $Y_s$ fields and contacting sides to
$\mathcal K_{I-P_W}$, then add the side through $J$. At a leaf the
$Y_s$ frame is $D_s=\mathcal K_{I-P_s}$.
\item Advance every side to $\mathcal K_I$ and subtract its $V$-row
sum of the full-frame $X$. Its endpoint is $\mathcal K_I Z_e$.
\end{enumerate}

To verify the word, denote the incoming logical $X$ values by $x$.
After the source sweep side $e$ holds
$\mathcal K_U(Z_e+(Vx)_e)$. Its contribution to a target $s$ after
the target sweep is $D_s(Z_e+(Vx)_e)$. The first subtraction cancels
its arbitrary incoming value. The remaining logical target increment
is $(RG+JV)x=x$ by \eqref{eq:binary-incidence}. Thus the raw main
endpoint is
\[
 \operatorname{diag}(D,D)
 \begin{pmatrix}I&0\\Q^{-1}&I\end{pmatrix},
 \qquad Q_s=\mathcal K_{P_s},\quad D_s=\mathcal K_{I-P_s}.
\]
Every side has its own full rank-eight endpoint. A side of dimensions
$(a,b)$ pays jumps $a,8-a-b,b$; rank zero is the identity. Each main
tree path pays $1,2,4$. This verifies arbitrary side inputs, not just
zero-initialized side fields.

The same chronology applies to the rational address permutations
$I+(R_3-I)\otimes P_U$, where
$R_3=\left(\begin{smallmatrix}0&-1\\1&-1\end{smallmatrix}\right)$.
Payload additions are still binary, so \eqref{eq:binary-incidence}
is unchanged. In inverse stages all frames, including complementary
and helper frames, are inverted. Replacing every coordinate slot by
$f$ columns preserves the projector relations and each common-frame
contact; it does not distribute a tensor power over a sum.

\subsection{Five stages and the partial endpoints}
For $a^3=1$, set $T_i=I+\lambda_i u_i\phi_i$ in the ring
$\mathbb Z[a]/(a^3-1)$, with chronological order
\[
\begin{array}{c|c|c|c}
i&u_i^T&\phi_i&\lambda_i\\\hline
1&(0,0,1,1)&(1,1,0,0)&a-1\\
2&(0,1,-1,-1)&(-1,-1,0,-1)&a^2-1\\
3&(0,1,-1,0)&(-1,0,0,-1)&a-1\\
4&(1,0,-1,0)&(-1,0,-1,0)&a^2-1\\
5&(1,0,0,0)&(0,0,-1,0)&a-1.
\end{array}
\]
Each $\phi_i u_i=0$, and coefficientwise multiplication gives
\begin{equation}
 T_5T_4T_3T_2T_1=\operatorname{diag}(a,a,a^2,a^2).
 \label{eq:five-bank-word}
\end{equation}
The three coefficients of every entry at each intermediate product
are supplied by \path{checks/expected/finite_certificate.json}.
They give an exact integer polynomial identity, rather than a test
on sampled eigenvalues.

Set $a=Q^{-1}$. A fixed bank basis turns the raw shear into each
listed direction; the free constant part of the shear is subtracted
by a binary bank update. Use five windows sharing the rank-one
projector $P$ and having disjoint rank-seven complements. Initially
apply each complement $D_i$ on all four banks. For stages four and
five instead use the unimodular bases
\[
 S_4=[-e_0,e_0-e_2,e_1,e_3],\qquad
 S_5=[-e_2,e_0,e_1,e_3]
\]
and apply $D_i$ on their two active banks only. Both transvections
preserve $E=\operatorname{diag}(1,0,1,0)$, so the endpoint is
\[
 D_1D_2D_3(I-E+ED_4D_5)\operatorname{diag}(a,a,a^2,a^2).
\]
The partial projector reflection on banks zero and one changes their
$a$ to $a^{-1}=Q$ and fixes every complement. Banks two and three
already have $a^2=Q$. Thus banks zero and two receive rank $36$
and banks one and three rank $22$; the other slots are identity.
The bank bases and the five window placements are fixed slot maps.
To match equal-rank phase jobs, index copies by the full finite
$\mathrm{GL}_{36}(\F_3)$ cover and use right multiplication for these
placements. This permutes copies and multiplies productive and paid
counts by the common fixed integer $C=|\mathrm{GL}_{36}(\F_3)|$.
For rational addresses take $C$ schematic copies with fixed right-route
symbols; their full residue cover is constructed after precompilation
in Appendix~\ref{app:movement}, without a finite rational matcher group.

There are five independent sides per rectangle, hence $18910$
rank-eight outputs, and $240$ main outputs at each of ranks $22,36$.
Sides contribute five jumps at each of $a,8-a-b,b$ per rectangle;
main flag paths contribute $1200$ calls at each of ranks $1,2,4$.
The partial-bank schedule leaves $720$ inactive rank-seven calls.
The $45$ owned increments contribute one call each at ranks one and
seven. These give the normalized profiles $T,P$ of Lemma~\ref{lem:unit}.
The actual covered program has the integer profiles $CT,CP$.
Multiplication by $C$ preserves the strict moment and rank thresholds;
it is included in the compiler's integer colors and finite constants.

\subsection{The strict moment certificate}
For a rational $x\in[1,2]$, put $y=(x-1)/(x+1)\le1/3$. Then
\[
 \log x=2\sum_{j=0}^{47}\frac{y^{2j+1}}{2j+1}+E,
 \qquad 0\le E\le\frac{18y^{97}}{8\cdot97}.
\]
Normalize each $36/r$ by powers of two and use these directed
bounds with binary scale $2^{48}$. For
$h=\log(36/r)/1300$ use
\[
 1+h+h^2/2\le e^h\le1+h+\frac{h^2}{2(1-h/3)}.
\]
Since $(r/36)^{s_0}=(r/36)e^h$, outward integer rounding gives the
two rational endpoints displayed in Lemma~\ref{lem:unit}. Their
numerator difference is $48274644696285>0$. The fixed checker prints
all per-rank intervals, so the strict moment is independently
reproducible with exact arithmetic. No growing-size running time is
inferred from this finite calculation.
